
Soit donc $P$ un torseur sous $V$, correspondant à une extension $\underline{F}$ de $\underline{E}$ par $\underline{O}_S$. L'algèbre $\underline{A}(P)$ de $P$ sur $\underline{O}_S$ est $\mathbb S(\underline{F})\otimes_{\mathbb S(\underline{O}_S)}\underline{O}_S$, soit $\mathbb S(\underline{F})/(s-1)\mathbb S(\underline{F})$, qui s'identifie à l'anneau des fractions homogènes de degré $0$ $\mathbb S(\underline{F})_{(s)}$. Par suite, $P$ s'identifie à l'ouvert $D_+(s)$ de $\mathbb P(\underline{F})$. Le complémentaire de $P$ dans $\mathbb P(\underline{F})$ est alors $\operatorname{Proj}(\mathbb S(\underline{F})/s\cdot \mathbb S(\underline{F}))$, qui s'identifie d'après (3.3.1) à $\mathbb P(\underline{E})$. On a donc les immersions :
\[
\begin{tikzcd}[column sep=large,row sep=small]
\mathbb P(\underline{E}) \arrow[r,hook] \arrow[dr] & \mathbb P(\underline{F}) \arrow[d] & P \arrow[l,hook'] \arrow[dl]\\
& S &
\end{tikzcd}
\]
où $\mathbb P(\underline{E})$ est fermé dans $\mathbb P(\underline{F})$, $P$ ouvert et complémentaire de $\mathbb P(\underline{E})$.

\textbf{Proposition 3.4.} \uline{Soient} $S$ \uline{un schéma noethérien}, $\underline{E}$ \uline{un} $\underline{O}_S$\uline{-Module localement libre de rang} $n$, $V=\mathbb V(\underline{E})$ \uline{le fibré vectoriel associé}, $P$ \uline{un torseur sous} $V$, $f$ \uline{le morphisme structural de} $P$. \uline{Alors l'homomorphisme canonique} :
\[
f^*:K_\bullet(S)\longrightarrow K_\bullet(P)
\]
\uline{est un isomorphisme}.

Avec les identifications de 3.3, on a la suite exacte :
\[
K_\bullet(\mathbb P(\underline{E}))\xrightarrow{i_*}K_\bullet(\mathbb P(\underline{F}))\longrightarrow K_\bullet(P)\longrightarrow 0,
\tag{3.4.1}
\]
où $i$ désigne l'immersion fermée $\mathbb P(\underline{E})\to \mathbb P(\underline{F})$ définie par $u$. D'après 3.2, on a $K_\bullet(\mathbb P(\underline{E}))=K_\bullet(S)\otimes_{K^{\bullet}(S)}K^{\bullet}(\mathbb P(\underline{E}))$, $K_\bullet(\mathbb P(\underline{F}))=K_\bullet(S)\otimes_{K^{\bullet}(S)}K^{\bullet}(\mathbb P(\underline{F}))$.
Comme $\mathbb P(\underline{E})$ et $\mathbb P(\underline{F})$ sont lisses, $i$ est une immersion régulière, et on a donc un homomorphisme $i_*:K^{\bullet}(\mathbb P(\underline{E}))\to K^{\bullet}(\mathbb P(\underline{F}))$; l'homomorphisme entre les $K_\bullet$ est alors obtenu en tensorisant l'homomorphisme entre les $K^{\bullet}$ par $K_\bullet(S)$ : il suffit pour le voir de montrer que pour $a\in K_\bullet(S)$ et
